\documentclass[11pt,a4paper]{article}
\usepackage[margin=1in]{geometry}
\usepackage{amsmath,amssymb,amsthm,bm,mathtools}
\usepackage{xeCJK}
\usepackage{tikz}
\usetikzlibrary{decorations.pathreplacing,arrows.meta}
\usepackage{hyperref}
\usepackage{enumitem}
\usepackage{booktabs}
\usepackage{xcolor}
\usepackage{tcolorbox}
%\setCJKmainfont{Noto Serif CJK SC}
\hypersetup{colorlinks=true,linkcolor=blue,citecolor=blue,urlcolor=blue}
\tcbset{colback=gray!4,colframe=gray!45,boxrule=0.5pt,arc=2mm}

\title{Akra--Bazzi 定理的一个简单而直观的证明\\[4pt]}
%\large 从“叶子数”、\(\log\)-size 坐标到 Stieltjes 积分，并与 Leighton 的归纳证明衔接}
\author{Dongbo Bu\\\texttt{dbu@ict.ac.cn}\\Institute of Computing Technology, \\Chinese Academy of Sciences}
\date{}

\newtheorem{theorem}{定理}
\newtheorem{remark}{注}
\newtheorem{proposition}{命题}

\begin{document}
\maketitle

\begin{abstract}
Akra-Bazzi是1998年正式发表（1996年就有一些初始版本）。1996年，MIT Tom Leighton写了一个Note，
给出一个适合本科生的证明；不过这个证明是``已知结论，然后证明"类型的；本文试图阐述如何一步一步``想"这个定理来。

核心是一句话：\textbf{总工作量 = 叶子数 \(\times\) 每个叶子沿 \(\log\)-size 递归深度累计的平均工作量。} 我们分三点描述：
\begin{enumerate}[label=(\arabic*)]
    \item 指数 \(p\) 是递归树\textbf{叶子数的增长指数}：规模为 \(x\) 的递归树有 \(\Theta(x^p)\) 个叶子；
    \item 递归按\textbf{比例}缩小问题，因此天然坐标不是 size 本身，而是 \(t=\log(\text{size})\)。在这个坐标中，几何级数式的递归层被“拉直”，从而离散求和才自然对应 Riemann 型积分；
    \item 一般 Akra--Bazzi 中不同 \(b_i\) 造成不规则的 \(\log\)-网格。严格地说，离散和先变成一个 Riemann--Stieltjes 积分，再在 renewal（更新）意义下连续化为普通积分。
\end{enumerate}
\end{abstract}

\section{定理的核心形式}
考虑递归式
\begin{equation}
T(x)=\sum_{i=1}^k a_i\,T(b_i x)+g(x),
\qquad a_i>0,\quad 0<b_i<1.
\end{equation}
Akra--Bazzi 的原始结果发表于 1998 年\cite{AkraBazzi1998}。设 \(p\) 是方程
\begin{equation}
\sum_{i=1}^k a_i b_i^p = 1
\end{equation}
的唯一实根。Akra--Bazzi 定理给出
\begin{equation}
T(x)=\Theta\!\left(x^p\left(1+\int_1^x \frac{g(u)}{u^{p+1}}\,du\right)\right).
\end{equation}

本文讨论最核心的无扰动形式。标准 Akra--Bazzi 定理允许递归参数中再带较小的加性扰动；这些技术细节不改变本文要解释的结构。

\section{第一步：\(p\) 就是叶子数的增长指数}
先忽略 \(g(x)\)，只看递归树的分支结构。记
\begin{equation}
L(x)=\text{规模为 }x\text{ 的结点最终产生的叶子数}.
\end{equation}
那么
\begin{equation}
L(x)=\sum_{i=1}^k a_i L(b_i x).
\end{equation}
若 \(L(x)\asymp x^p\)，代入得到
\begin{equation}
x^p \asymp \sum_{i=1}^k a_i(b_i x)^p
=x^p\sum_{i=1}^k a_i b_i^p.
\end{equation}
因此恰好需要
\begin{equation}
\sum_{i=1}^k a_i b_i^p=1.
\end{equation}

\begin{tcolorbox}
\textbf{解释：}Akra--Bazzi 中的 \(p\) 不是为了配平公式而引入的神秘指数；它刻画递归树叶子数随输入规模的增长速度：
\[
\#\text{leaves}=\Theta(x^p).
\]
\end{tcolorbox}

\section{教学热身：所有缩放因子相同时，离散和怎样变成积分}
先考虑普通 Master-Theorem 形式
\begin{equation}
T(x)=aT(x/b)+g(x),\qquad b>1.
\end{equation}
此时
\begin{equation}
p=\log_b a,
\qquad ab^{-p}=1.
\end{equation}
定义
\begin{equation}
F(x)=\frac{T(x)}{x^p},
\qquad
h(x)=\frac{g(x)}{x^p}.
\end{equation}
则
\begin{equation}
F(x)=F(x/b)+h(x).
\end{equation}
若令 \(x=b^m\)，反复展开得到精确离散式
\begin{equation}
F(b^m)=F(1)+\sum_{j=1}^{m} h(b^j).
\end{equation}

\subsection{为什么不能直接在 size 轴上把它看成普通 Riemann 和？}
在原始 size 轴上，采样点是
\begin{equation}
1,b,b^2,\ldots,b^m=x.
\end{equation}
相邻点之间的宽度
\begin{equation}
\Delta u_j=b^j-b^{j-1}=(b-1)b^{j-1}
\end{equation}
随 \(u\) 成比例增长。因此这些点在 size 轴上极不均匀。

严格说，Riemann 积分本身\emph{并不要求等距划分}；但是一个\textbf{没有显式区间宽度因子的裸求和}
\begin{equation}
\sum_j f(u_j)
\end{equation}
若想直接解释成 Riemann sum，就需要各采样点代表近似相同的“坐标宽度”。而在 size 轴上，递归层显然没有这个性质。

\begin{tcolorbox}
\textbf{关键问题不是“积分必须等距”，而是：}\par
递归式给我们的离散和对每一层的权重都是 1。要把这种“等权离散和”解释成 Riemann 型积分，需要找到一个坐标，使相邻递归层对应的坐标宽度相同或至少具有稳定的平均密度。
\end{tcolorbox}

\subsection{为什么取 \(\log\)：把乘法递归拉直}
递归每一步保持的是\textbf{比例}：
\begin{equation}
u\longmapsto u/b.
\end{equation}
所以递归真正自然的“距离”不是两个规模之差，而是规模之比。取
\begin{equation}
t=\log u
\end{equation}
正好把“比”变成“差”：
\begin{equation}
\log u-\log(u/b)=\log b.
\end{equation}
于是
\begin{equation}
1,b,b^2,\ldots,b^m
\end{equation}
在 \(\log\)-size 轴上变成
\begin{equation}
0,\lambda,2\lambda,\ldots,m\lambda,
\qquad \lambda=\log b,
\end{equation}
即严格等距网格。

\begin{center}
\begin{tikzpicture}[x=0.82cm,y=0.9cm]
  \node[anchor=east,font=\small\bfseries] at (-0.15,1.7) {size 轴：};
  \draw[-{Latex[length=2mm]},thick] (0,1.7)--(13.0,1.7) node[right] {$u$};
  \foreach \x/\lab in {0.6/$1$,1.3/$b$,2.8/$b^2$,5.8/$b^3$,11.8/$b^4$}{
    \filldraw[fill=orange!80!black,draw=orange!80!black] (\x,1.7) circle (1.8pt);
    \node[above] at (\x,1.7) {\lab};
  }
  \draw[decorate,decoration={brace,amplitude=4pt},gray] (0.6,1.1)--(1.3,1.1);
  \draw[decorate,decoration={brace,amplitude=4pt},gray] (1.3,1.1)--(2.8,1.1);
  \draw[decorate,decoration={brace,amplitude=4pt},gray] (2.8,1.1)--(5.8,1.1);
  \draw[decorate,decoration={brace,amplitude=4pt},gray] (5.8,1.1)--(11.8,1.1);
  \node[font=\footnotesize,gray] at (6.2,0.55) {相邻层间距随规模增大：不是等权 Riemann 网格};

  \node[anchor=east,font=\small\bfseries] at (-0.15,-0.55) {$\log$-size 轴：};
  \draw[-{Latex[length=2mm]},thick] (0,-0.55)--(13.0,-0.55) node[right] {$t=\log u$};
  \foreach \x/\lab in {1.0/$0$,3.5/$\lambda$,6.0/$2\lambda$,8.5/$3\lambda$,11.0/$4\lambda$}{
    \filldraw[fill=blue!70!black,draw=blue!70!black] (\x,-0.55) circle (1.8pt);
    \node[above] at (\x,-0.55) {\lab};
  }
  \foreach \a/\b in {1.0/3.5,3.5/6.0,6.0/8.5,8.5/11.0}{
    \draw[<->,gray] (\a,-1.12)--(\b,-1.12);
  }
  \node[font=\footnotesize,gray] at (6.0,-1.55) {每个递归层对应同样的坐标宽度 $\lambda=\log b$};
\end{tikzpicture}
\end{center}

现在离散和
\begin{equation}
\sum_{j=1}^{m} h(e^{j\lambda})
\end{equation}
才是标准的等距网格求和。乘上网格宽度 \(\lambda\) 后，得到 Riemann 和：
\begin{equation}
\lambda\sum_{j=1}^{m} h(e^{j\lambda})
\asymp
\int_0^{t}h(e^s)\,ds,
\qquad t=m\lambda.
\end{equation}
所以
\begin{equation}
F(e^t)
=\Theta\!\left(1+\frac{1}{\lambda}\int_0^t h(e^s)\,ds\right).
\end{equation}
令 \(u=e^s\)，于是
\begin{equation}
ds=\frac{du}{u},
\end{equation}
得到
\begin{equation}
F(x)=\Theta\!\left(1+\int_1^x\frac{g(u)}{u^{p+1}}\,du\right).
\end{equation}

\begin{tcolorbox}
\textbf{这一节最重要的一句话：}\par
\textbf{取 \(\log\) 不是为了“凑出” \(1/u\)，而是为了把乘法递归的几何级数网格拉直成近似均匀的加法网格，使等权离散求和可以自然解释为 Riemann 型积分。}
\end{tcolorbox}

\section{一般情形：先除以叶子数}
回到
\begin{equation}
T(x)=\sum_{i=1}^k a_i T(b_i x)+g(x).
\end{equation}
定义
\begin{equation}
F(x)=\frac{T(x)}{x^p},
\qquad
q_i=a_i b_i^p,
\qquad
h(x)=\frac{g(x)}{x^p}.
\end{equation}
因为 \(\sum_i a_i b_i^p=1\)，有
\begin{equation}
\sum_{i=1}^k q_i=1.
\end{equation}
递归式变成
\begin{equation}
F(x)=\sum_{i=1}^k q_i F(b_i x)+h(x).
\end{equation}
这里 \(q_i\) 可以解释为第 \(i\) 类孩子最终贡献的叶子比例。因此可以把它看成：随机选取一个最终叶子时，下一步走向第 \(i\) 类孩子的概率。

\section{随机叶子路径：得到精确的离散求和}
令 \(I_1,I_2,\ldots\) 独立同分布，满足
\begin{equation}
\Pr(I_j=i)=q_i.
\end{equation}
沿一条随机叶子路径，问题规模依次为
\begin{equation}
X_0=x,
\qquad
X_1=b_{I_1}x,
\qquad
X_2=b_{I_2}b_{I_1}x,
\quad\ldots
\end{equation}
直到规模降到常数范围。反复展开归一化递归式，可得
\begin{equation}
F(x)
=
B(x)
+
\mathbb{E}\!\left[\sum_{j<\tau_x} h(X_j)\right],
\end{equation}
其中 \(\tau_x\) 是到达边界规模的停止时间，\(B(x)=\Theta(1)\) 是边界项。

因此，在任何积分出现以前，我们首先得到的是一个\textbf{完全离散}的表达：
\begin{quote}
\textbf{\(T(x)/x^p\) = 一条随机叶子路径上，各级“单位叶子工作” \(h\) 的离散累加。}
\end{quote}

\section{一般 Akra--Bazzi 中，\(\log\)-size 仍然是自然坐标}
令
\begin{equation}
t=\log x,
\qquad
\lambda_i=\log(1/b_i)>0.
\end{equation}
一次递归
\begin{equation}
x\mapsto b_i x
\end{equation}
变成
\begin{equation}
t\mapsto t-\lambda_i.
\end{equation}
定义
\begin{equation}
S_0=0,
\qquad
S_j=\lambda_{I_1}+\cdots+\lambda_{I_j}.
\end{equation}
则
\begin{equation}
X_j=e^{t-S_j}.
\end{equation}
所以
\begin{equation}
F(e^t)
=
B(e^t)
+
\mathbb{E}\!\left[
\sum_{S_j<t} h\!\left(e^{t-S_j}\right)
\right].
\end{equation}

与等 \(b\) 情形相比，这里的 \(\log\)-网格不再严格等距，因为每一步的长度从有限集合
\begin{equation}
\{\lambda_1,\ldots,\lambda_k\}
\end{equation}
中选择。但这仍比原始 size 轴自然得多：乘法路径已经变成了\textbf{加法随机游走}。

平均步长为
\begin{equation}
\mu
=\sum_{i=1}^k q_i\lambda_i
=\sum_{i=1}^k a_i b_i^p\log(1/b_i).
\end{equation}
因此从宏观尺度看，\(\log\)-size 轴上递归层的平均密度约为 \(1/\mu\)。这正是一般情形从“近似等距网格”走向积分所需的替代条件。

\begin{center}
\begin{tikzpicture}[x=0.72cm,y=0.9cm]
    \node[anchor=east,font=\small\bfseries] at (-0.1,1.8) {等 $b$：};
    \draw[-{Latex[length=2mm]},thick] (0,1.8) -- (13.4,1.8) node[right] {$t=\log(\text{size})$};
    \foreach \x in {1.0,3.0,5.0,7.0,9.0,11.0,13.0} {
      \filldraw[fill=blue!70!black,draw=blue!70!black] (\x,1.8) circle (1.7pt);
    }
    \draw[<->,thick,gray] (9.0,1.18)--(11.0,1.18);
    \node[below,gray] at (10.0,1.18) {$\lambda$};
    \node[above,font=\footnotesize] at (7.0,2.0) {严格等距：普通 Riemann 网格};

    \node[anchor=east,font=\small\bfseries] at (-0.1,-0.3) {一般 AB：};
    \draw[-{Latex[length=2mm]},thick] (0,-0.3) -- (13.4,-0.3) node[right] {$t=\log(\text{size})$};
    \foreach \x in {0.8,2.0,4.3,5.5,7.8,9.0,10.2,12.5} {
      \filldraw[fill=orange!85!black,draw=orange!85!black] (\x,-0.3) circle (1.7pt);
    }
    \draw[<->,thick,gray] (9.0,-0.92)--(10.2,-0.92);
    \node[below,gray] at (9.6,-0.92) {$\lambda_{i_j}$};
    \node[above,font=\footnotesize] at (7.0,-0.1) {不规则，但长期具有稳定平均密度 $1/\mu$};
\end{tikzpicture}
\end{center}

\section{正文中的桥梁：离散和先变成 Stieltjes 积分，再变成普通积分}
现在不能把离散和直接写成普通积分。最自然的中间步骤是引入“递归层计数函数”。

对一条给定路径，定义
\begin{equation}
N(s)=\#\{j\ge 0:S_j\le s\}.
\end{equation}
它是一个阶梯函数，每经过一个递归层位置 \(S_j\) 就跳 1。于是有精确等式
\begin{equation}
\sum_{S_j<t} h\!\left(e^{t-S_j}\right)
=
\int_{[0,t)} h\!\left(e^{t-s}\right)\,dN(s).
\end{equation}
右边就是 Riemann--Stieltjes 积分。这里没有连续化，也没有近似：\textbf{阶梯函数的每一次跳跃，恰好记录一个离散求和项。}

对随机路径取期望，定义 renewal function
\begin{equation}
U(s)=\mathbb{E}[N(s)].
\end{equation}
于是
\begin{equation}
\mathbb{E}\!\left[
\sum_{S_j<t} h\!\left(e^{t-S_j}\right)
\right]
=
\int_{[0,t)} h\!\left(e^{t-s}\right)\,dU(s).
\end{equation}
这一步仍然是精确的。

更新理论告诉我们：长期看，renewal points 在 \(\log\)-size 轴上的平均密度是 \(1/\mu\)。最基本的形式是
\begin{equation}
\frac{U(t)}{t}\longrightarrow \frac{1}{\mu}.
\end{equation}
在 Akra--Bazzi 所需的正则性条件下，这种稳定平均密度足以把上面的 Stieltjes 积分在 \(\Theta\) 意义下替换为
\begin{equation}
\frac{1}{\mu}\int_0^t h(e^s)\,ds.
\end{equation}
若追求更精细的极限常数，则需要相应的 renewal theorem / key renewal theorem 条件；格点情形还可能带有有界周期振荡，但不影响 \(\Theta\) 结论。

因此
\begin{equation}
F(e^t)
=
\Theta\!\left(1+\int_0^t h(e^s)\,ds\right).
\end{equation}
代入 \(h(u)=g(u)/u^p\)，令 \(u=e^s\)，则 \(ds=du/u\)：
\begin{align}
F(x)
&=\Theta\!\left(1+\int_0^{\log x}\frac{g(e^s)}{e^{ps}}\,ds\right)\\
&=\Theta\!\left(1+\int_1^x\frac{g(u)}{u^p}\frac{du}{u}\right)\\
&=\Theta\!\left(1+\int_1^x\frac{g(u)}{u^{p+1}}\,du\right).
\end{align}
最后乘回 \(x^p\)：
\begin{equation}
\boxed{
T(x)=\Theta\!\left(
 x^p\left[1+\int_1^x\frac{g(u)}{u^{p+1}}\,du\right]
\right).
}
\end{equation}

\section{公式中每一部分从哪里来？}
\begin{center}
\begin{tabular}{lll}
\toprule
公式中的部分 & 来源 & 含义\\
\midrule
\(x^p\) & 叶子数 & 递归树最终有多少个叶子\\
\(g(u)/u^p\) & 叶子归一化 & 规模为 \(u\) 时，每个叶子分摊的本层工作\\
\(d(\log u)\) & 自然递归坐标 & 递归层按比例缩放，而不是按差值缩放\\
\(du/u\) & 变量代换 & \(d(\log u)=du/u\)\\
\(1/\mu\)（隐藏常数） & 层密度 & 单位 \(\log\)-size 长度中的平均递归层数\\
\bottomrule
\end{tabular}
\end{center}

所以 Akra--Bazzi 的积分，更自然的原始写法其实是
\begin{equation}
\int \frac{g(u)}{u^p}\,d(\log u),
\end{equation}
标准形式只是把 \(d(\log u)=du/u\) 代进去。

\section{一个最简单的检验：临界情形}
在等 \(b\) 情形，若
\begin{equation}
g(x)=x^p,
\end{equation}
则 \(h(x)=1\)。每个递归层给每个叶子贡献同样的常数工作。递归深度为
\begin{equation}
\log_b x=\frac{\log x}{\log b}.
\end{equation}
所以
\begin{equation}
T(x)=\Theta(x^p\log x).
\end{equation}
积分公式也给出
\begin{equation}
x^p\int_1^x\frac{u^p}{u^{p+1}}\,du
=x^p\log x.
\end{equation}
完全吻合。


\section{与 Leighton（1996）的归纳验证证明如何衔接}
上面的路线是在回答：\emph{Akra--Bazzi 公式为什么会长成这个样子？}
Leighton 在 1996 年的讲义中给出了另一条非常简洁的路线：\emph{先接受这个候选形式，再直接验证它确实满足原递归的增长规律}\cite{Leighton1996}。

令
\begin{equation}
I(x)=\int_1^x \frac{g(u)}{u^{p+1}}\,du,
\qquad
\Phi(x)=x^p\,[1+I(x)].
\end{equation}
Leighton 证明中的关键引理是：在其 polynomial-growth 条件下，对每个固定的 \(i\)，存在与 \(x\) 无关的正常数，使得
\begin{equation}
\boxed{
 x^p\int_{b_i x}^{x}\frac{g(u)}{u^{p+1}}\,du
 =\Theta(g(x)).
}
\end{equation}
这条引理说明：从规模 \(x\) 走一步到 \(b_i x\) 时，候选函数的“积分部分”恰好损失一个 \(\Theta(g(x))\) 的量级。

事实上，记 \(q_i=a_i b_i^p\)，则 \(\sum_i q_i=1\)。于是
\begin{align}
\sum_{i=1}^k a_i\Phi(b_i x)
&=x^p\sum_{i=1}^k q_i\,[1+I(b_i x)]\\
&=x^p[1+I(x)]
  -x^p\sum_{i=1}^k q_i
      \int_{b_i x}^{x}\frac{g(u)}{u^{p+1}}\,du.
\end{align}
因此
\begin{equation}
\boxed{
\Phi(x)-\sum_{i=1}^k a_i\Phi(b_i x)=\Theta(g(x)).
}
\end{equation}
也就是说，\(\Phi(x)\) 自己就近似满足
\begin{equation}
T(x)=\sum_i a_iT(b_i x)+g(x),
\end{equation}
于是通过上下界归纳即可得到 \(T(x)=\Theta(\Phi(x))\)。这就是 Leighton 证明的核心。

\subsection*{为什么这个引理与本文的推导正好吻合？}
本文从递归树出发，把
\begin{equation}
\frac{T(x)}{x^p}
\end{equation}
解释成沿随机叶子路径累计的“单位叶子工作”。在 \(\log\)-size 坐标中，从 \(x\) 到 \(b_i x\) 正好跨越一个固定宽度
\begin{equation}
\log x-\log(b_i x)=\log(1/b_i).
\end{equation}
因此，一次递归步所跨越的这段 \(\log\)-尺度中，累计的单位叶子工作应该与当前层的 \(g(x)/x^p\) 同阶。乘回 \(x^p\)，就正是 Leighton 的局部积分引理。

换句话说，这个引理不仅是归纳证明的技术工具，也可以看成对前面“叶子数 + \(\log\)-size”推导的一个\textbf{local consistency check}：
\begin{quote}
从 \(x\) 缩到 \(b_i x\) 的一段自然递归尺度，所累计的工作，应该与当前这一层的非递归工作 \(g(x)\) 同阶。
\end{quote}

\begin{tcolorbox}
\textbf{两条路线的分工：}
\[
\begin{array}{ccl}
\text{本文路线} &:& \text{recursion tree}\to\text{leaves}\to\log\text{-size}\to\text{integral}
\\[1mm]
&&\text{回答“公式是怎样被发现的？”}\\[2mm]
\text{Leighton 1996} &:& \text{candidate }\Phi\to\text{local integral lemma}\to\text{induction}
\\[1mm]
&&\text{回答“怎样简洁地验证公式是对的？”}
\end{array}
\]
因此，两者最适合的组合是：\textbf{先推导（discovery），再归纳验证（verification）}。
\end{tcolorbox}

\section{总结：五步重新“发明”Akra--Bazzi}
\begin{enumerate}[label=\arabic*.]
    \item \textbf{数叶子：}\(\#\text{leaves}=\Theta(x^p)\)，其中 \(p\) 满足 \(\sum_i a_i b_i^p=1\)；
    \item \textbf{除以叶子数：}\(F(x)=T(x)/x^p\)，把问题变成“每个叶子的平均总代价”；
    \item \textbf{沿随机叶子路径离散求和：}\(F\) 等于各级单位叶子工作 \(h\) 的累计；
    \item \textbf{取 \(\log\)-size：}把乘法递归拉直，使等权离散层成为等距或具有稳定平均密度的网格；
    \item \textbf{离散到连续：}等 \(b\) 时直接得到 Riemann 和；一般情形先精确写成 Stieltjes 积分，再利用 renewal density 得到普通积分。
\end{enumerate}

\begin{tcolorbox}
\centering
\textbf{最简洁的理解：}\par
\vspace{2mm}
\(p\) 来自\textbf{叶子数}；\qquad
\(\log\) 来自\textbf{把乘法递归拉直}；\qquad
\(1/u\) 来自\textbf{把 \(d\log u\) 换回 \(du\)}。\par
\vspace{1mm}
积分不是突然出现的：它是\textbf{离散递归层在自然坐标中的连续化}。
\end{tcolorbox}

\appendix
\section{附录：什么是 Riemann--Stieltjes 积分？}
\subsection{从普通 Riemann 积分出发}
普通 Riemann 积分把区间 \([a,b]\) 划分为
\begin{equation}
a=x_0<x_1<\cdots<x_n=b,
\end{equation}
在每个小区间取一点 \(\xi_j\in[x_{j-1},x_j]\)，形成和
\begin{equation}
\sum_{j=1}^n f(\xi_j)\,(x_j-x_{j-1}).
\end{equation}
当划分的最大宽度趋于 0 时，该和趋向
\begin{equation}
\int_a^b f(x)\,dx.
\end{equation}
这里累加的“质量”来自横轴长度 \(\Delta x\)。

\subsection{Stieltjes 积分：把 \(\Delta x\) 换成 \(\Delta A\)}
给定一个单调不减函数 \(A(x)\)，Riemann--Stieltjes 和定义为
\begin{equation}
\sum_{j=1}^n f(\xi_j)\,[A(x_j)-A(x_{j-1})].
\end{equation}
在适当条件下，其极限记为
\begin{equation}
\int_a^b f(x)\,dA(x).
\end{equation}
所以它的直观意义是：
\begin{quote}
普通 Riemann 积分按“横轴长度”累加；Stieltjes 积分按“\(A\) 增加了多少”累加。
\end{quote}
如果 \(A(x)=x\)，就退化为普通 Riemann 积分。

\subsection{阶梯函数时，Stieltjes 积分就是离散求和}
假设有离散点
\begin{equation}
s_1<s_2<s_3<\cdots
\end{equation}
并定义计数函数
\begin{equation}
N(s)=\#\{j:s_j\le s\}.
\end{equation}
则 \(N\) 是阶梯函数，在每个 \(s_j\) 处跳 1。于是
\begin{equation}
\int f(s)\,dN(s)=\sum_j f(s_j).
\end{equation}
如果某一点的跳跃大小是 \(c_j\)，则该点贡献 \(c_jf(s_j)\)。

这正是本文需要 Stieltjes 积分的原因：它允许我们把“对递归层位置逐点求和”写成一个\textbf{完全等价的积分形式}，而没有提前做任何连续近似。

\subsection{在 Akra--Bazzi 证明中的角色}
沿一条随机叶子路径，\(\log\)-size 轴上的递归层位置为
\begin{equation}
0=S_0<S_1<S_2<\cdots.
\end{equation}
定义
\begin{equation}
N(s)=\#\{j:S_j\le s\}.
\end{equation}
则
\begin{equation}
\sum_{S_j<t} h\!\left(e^{t-S_j}\right)
=
\int_{[0,t)} h\!\left(e^{t-s}\right)\,dN(s).
\end{equation}
取期望并令 \(U(s)=\mathbb{E}N(s)\)，得到
\begin{equation}
\mathbb{E}\sum_{S_j<t} h\!\left(e^{t-S_j}\right)
=
\int_{[0,t)} h\!\left(e^{t-s}\right)\,dU(s).
\end{equation}
此时，renewal theory 才负责告诉我们 \(U\) 在大尺度上的平均斜率约为 \(1/\mu\)，从而把 \(dU(s)\) 在宏观上替换为 \((1/\mu)ds\)。

因此，整个“离散到连续”的逻辑链是
\begin{equation}
\boxed{
\text{离散求和}
\;\longrightarrow\;
\text{Stieltjes 积分（精确）}
\;\longrightarrow\;
\text{普通积分（连续化）}.
}
\end{equation}

\subsection{一个小心点：不要把 \(U(t)/t\to1/\mu\) 直接当成逐点导数}
从
\begin{equation}
\frac{U(t)}{t}\to\frac{1}{\mu}
\end{equation}
并不能形式上直接推出 \(U'(t)\to1/\mu\)，因为 \(U\) 甚至可能不是光滑函数。正确理解是：在足够大的区间上，renewal measure 的\textbf{平均质量密度}趋于 \(1/\mu\)。在本文只追踪 \(\Theta\) 量级且 \(h\) 满足相应正则性时，这已足够；若要得到精确渐近常数，则应使用更精细的 renewal theorem。


\renewcommand{\refname}{参考文献}
\begin{thebibliography}{9}

\bibitem{AkraBazzi1998}
Mohamad Akra and Louay Bazzi.
\newblock ``On the Solution of Linear Recurrence Equations.''
\newblock \emph{Computational Optimization and Applications}, 10(2):195--210, 1998.
\newblock DOI: \href{https://doi.org/10.1023/A:1018373005182}{10.1023/A:1018373005182}.

\bibitem{Leighton1996}
Tom Leighton.
\newblock ``Notes on Better Master Theorems for Divide-and-Conquer Recurrences.''
\newblock Mathematics Department and Laboratory for Computer Science, Massachusetts Institute of Technology, October 9, 1996.
\newblock \url{https://courses.csail.mit.edu/6.046/spring04/handouts/akrabazzi.pdf}.

\end{thebibliography}

\end{document}
